% Part: first-order-logic
% Chapter: tableaux
% Section: provability-quantifiers

\documentclass[../../../include/open-logic-section]{subfiles}

\begin{document}

\olfileid{fol}{tab}{qpr}
\olsection{\usetoken{S}{derivability} અને પરિમાણકો}

\begin{explain}
  પૂર્ણતાના પ્રમેય માટે એ પણ જરૂરી છે કે ટેબ્લોના નિયમો આ વિભાગમાં
  ~$\Proves$ વિશે સ્થાપિત કરેલી હકીકતો આપે.
\end{explain}

\begin{thm}
\ollabel{thm:strong-generalization} જો $c$ એવો અચળ હોય જે $\Gamma$ કે
~$!A(x)$માં આવતો નથી અને $\Gamma \Proves !A(c)$ હોય, તો
$\Gamma \Proves \lforall[x][!A(x)]$.
\end{thm}

\begin{proof}
ધારો કે $\Gamma \Proves !A(c)$, એટલે કે એવાં $!B_1$, \dots,
$!B_n \in \Gamma$ અને નીચેના ગણ માટે બંધ !!{tableau} છે:
\begin{align*}
\{ \sFmla{\False}{!A(c)}, &
\sFmla{\True}{!B_1}, \dots, \sFmla{\True}{!B_n} \}.
\intertext{આપણે દર્શાવવાનું છે કે નીચેના ગણ માટે પણ બંધ !!{tableau} છે:}
\{ \sFmla{\False}{\lforall[x][!A(x)]}, &
\sFmla{\True}{!B_1}, \dots, \sFmla{\True}{!B_n} \}.
\end{align*}
બંધ !!{tableau} લો, તેની પહેલી ધારણાને
$\sFmla{\False}{\lforall[x][!A(x)]}$થી બદલો અને ધારણાઓ પછી
$\sFmla{\False}{!A(c)}$ દાખલ કરો.
\begin{center}
\begin{tableau}{not line numbering}
  [\sFmla{\False}{\formula{A}(c)}
    [\sFmla{\True}{\formula{B}_1}
      [\vdots
      [\sFmla{\True}{\formula{B}_n},
        [][][]]]]]
\end{tableau}\qquad
\begin{tableau}{not line numbering}
  [\sFmla{\False}{\lforall[x][\formula{A}(x)]}
    [\sFmla{\True}{\formula{B}_1}
      [\vdots
        [\sFmla{\True}{\formula{B}_n},
          [\sFmla{\False}{\formula{A}(c)}
        [][][]]]]]]
\end{tableau}
\end{center}
ટેબ્લો હજી બંધ છે, કારણ કે પહેલાં ધારણાઓ તરીકે ઉપલબ્ધ બધાં
!!{sentence}s હજી !!{tableau}ની ટોચે ઉપલબ્ધ છે. દાખલ કરેલી પંક્તિ
$\TRule{\False}{\lforall}$ના યોગ્ય ઉપયોગનું પરિણામ છે, કારણ કે
!!{constant}~$c$ એ $!B_1$, \dots, $!B_n$ કે
$\lforall[x][!A(x)]$માં આવતો નથી; એટલે કે નવા !!{tableau}માં દાખલ
કરેલી પંક્તિની ઉપર આવતો નથી.
\end{proof}

\begin{prop}
\ollabel{prop:provability-quantifiers}
\begin{tagenumerate}{prvEx,prvAll}
\tagitem{prvEx}{$!A(t) \Proves \lexists[x][!A(x)]$.}{}
\tagitem{prvAll}{$\lforall[x][!A(x)] \Proves !A(t)$.}{}
\end{tagenumerate}
\end{prop}

\begin{proof}
\begin{tagenumerate}{prvEx,prvAll}
\tagitem{prvEx}{$\sFmla{\False}{\lexists[x][!A(x)]},
  \sFmla{\True}{!A(t)}$ માટેનો બંધ !!{tableau} આ છે:
  \begin{oltableau}
    [\sFmla{\False}{\lexists[x][\formula{A}(x)]}, just = \TAss
      [\sFmla{\True}{\formula{A}(t)}, just = \TAss
        [\sFmla{\False}{\formula{A}(t)}, just = {\TRule{\False}{\lexists}[1]}, close]]]
    \end{oltableau}}{}
\tagitem{prvAll}{$\sFmla{\False}{\formula{A}(t)},
  \sFmla{\True}{\lforall[x][!A(x)]}$ માટેનો બંધ !!{tableau} આ છે:
  \begin{oltableau}
    [\sFmla{\False}{\formula{A}(t)}, just = \TAss
      [\sFmla{\True}{\lforall[x][\formula{A}(x)]}, just = \TAss
        [\sFmla{\True}{\formula{A}(t)}, just = {\TRule{\True}{\lforall}[2]}, close]]]
    \end{oltableau}}{}
\end{tagenumerate}
\sourcecorrection{OLTAB-008}{મૂળની
\texttt{provability-quantifiers.tex}, પંક્તિ~79માં બીજા ચિહ્નિત સૂત્ર
પછીનું અલ્પવિરામ ગણની ગણિતીય સીમાની અંદર મૂકાયું છે. અહીં એ વધારાનું
અલ્પવિરામ દૂર કર્યું છે.}
\end{proof}

\end{document}
